% UNREGISTERED A.50; full independent root audit required before inclusion.
\subsection{The degree-twenty-seven Heisenberg cover}

\begin{statement}{A.50}{Proved}{All normalizers in an explicit genus-ten cover}
Let \(C\) be the smooth proper normalization of
\[
 y^2=x^3-x,\qquad z^3=\psi_3(x):=3x^4-6x^2-1,
 \qquad q:C\longrightarrow E.                         \tag{A.50.1}
\]
Put \(S=q^{-1}(E[3])\) and \(X=C\setminus S\).
Then \([3]q:X\to U\) is a connected geometric Galois
cover with Heisenberg deck group of order \(27\),
proper genus \(10\), and nine cusps.
For every coefficient prime, the nontrivial closed
procyclic subgroups of its geometric pro-prime group
with open arithmetic normalizer are precisely the
nontrivial closed subgroups of cusp inertia conjugates.
Its full pro-three fixed subgroup under every indicated
open derived group and cusp packet is the selected
cusp inertia, by A.22.

The proper Jacobian has an isogeny decomposition
\[
              J_C\sim E\times E_0^3\times B^3,
       \qquad E_0:3s^2=z^3+4,\quad \dim B=2,           \tag{A.50.2}
\]
after finite number-field extension. The surface
\(B\) has an action by \(\mathbf Q(\sqrt{-3})\)
and has no geometric CM abelian subvariety.
The \(E_0^3\) part is detected by actual punctured
maps, not merely by correspondences of Jacobians.
The curve also has potentially good reduction at two
with special two-rank zero.
No fixed-subgroup assertion at coefficient primes
other than three is inferred here from the normalizer
calculation.
\end{statement}

\begin{proof}
\paragraph{The actual cover and its deck group.}
The division polynomial has divisor
\[
 \operatorname{div}_E(\psi_3)
       =\sum_{T\in E[3]\setminus\{O\}}(T)-8(O).
                                                               \tag{A.50.3}
\]
Its eight finite zeroes and its pole have orders prime
to three. Thus \(q\) is connected of degree three,
totally ramified at precisely the nine points of
\(E[3]\). Each has a unique preimage. Riemann--Hurwitz
therefore gives \(2g(C)-2=9(3-1)=18\).
The composite with the etale degree-nine isogeny
\([3]\) ramifies only above \(O\). Removing its
full fibre gives the asserted degree and nine cusps.

Here is an explicit construction of all the deck lifts.
Let \(T=(a,b)\in E[3]\setminus\{O\}\), so
\(3a^4-6a^2-1=0\). The quotient isogeny with kernel
\(\langle T\rangle\) has coordinates
\[
 \begin{split}
 X_a&=\frac{2(x^3-x+a^3-a)}{(x-a)^2}-x-4a,\\
 Y_a&=y\frac{dX_a}{dx},\\
 E_a&:Y_a^2=X_a^3+(9-30a^2)X_a+42a-70a^3.
 \end{split}                                                   \tag{A.50.4}
\]
The first expression is the sum-of-translates formula
\(x(P)+x(P+T)+x(P-T)-2x(T)\); hence it is invariant
under translation by \(T\). Direct substitution gives
\[
               \psi_3(x)=3(x-a)^3(X_a+3a).          \tag{A.50.5}
\]
Consequently translation by \(T\) lifts to
\[
       \widetilde T(P,z)=
       \left(P+T,\;z\frac{x(P+T)-a}{x(P)-a}\right).
                                                               \tag{A.50.6}
\]
The product of the three displayed ratios telescopes,
so \(\widetilde T^3=1\). This formula is a rational
function-field automorphism and therefore extends to
\(C\). It commutes with the central automorphism
\(c:z\mapsto\zeta_3z\).

Choose \(T\ne O\) and put \(T'=iT\). The two
points form a basis of \(E[3]\), since \(i\) has
irreducible polynomial \(t^2+1\) over \(\mathbf F_3\).
Their lifts and \(c\) give all \(27\) deck
automorphisms of \([3]q\): their images contain all
nine translations, and their central kernel contains
\(\mu_3\). This group cannot be abelian. Otherwise
the inertia about the sole branch point on \(E\)
would be trivial, since that inertia is a commutator
in the once-punctured elliptic fundamental group.
It instead has order three by (A.50.3).
The two lifts have order three, their commutator is
central of order three, and the usual class-two
power formula shows every element has order dividing
three. This is precisely the Heisenberg group
\(H_3\) of order \(27\).
All these maps are defined over
\(L=\mathbf Q(i,E[3])\), which contains \(\zeta_3\)
by the Weil pairing; enlarging \(L\) later is harmless.
A.22 now gives the stated full pro-three fixed subgroup.

\paragraph{An actual once-punctured CM target.}
The order-four automorphism
\[
               \tau(x,y,z)=(-x,iy,z)
                                                               \tag{A.50.7}
\]
has quotient map
\[
          q_0:C\longrightarrow E_0,
          \qquad (x,y,z)\longmapsto(z,s=x^2-1),
          \qquad 3s^2=z^3+4.                     \tag{A.50.8}
\]
The identity is immediate from (A.50.1). Generically
there are four choices of \((x,y)\), permuted
transitively by \(\tau\), so this is the exact
quotient of degree four. With \(O\) its point at
infinity, \(E_0\) has full CM by
\(\mathbf Z[\zeta_3]\); equivalently it is
\(Y^2=X^3+108\) under \(X=3z,Y=9s\).
Its automorphism \([\zeta_3]\) is
\((z,s)\mapsto(\zeta_3z,s)\).

The three points
\[
          O,\quad (0,2/\sqrt3),\quad(0,-2/\sqrt3)
\]
are exactly \(\ker[1-\zeta_3]\). The two finite
points have four preimages each, the eight zeroes in
(A.50.3); infinity has the unique remaining cusp,
with ramification index four. Hence
\[
      [1-\zeta_3]\circ q_0:X\longrightarrow E_0\setminus\{O\}
                                                               \tag{A.50.9}
\]
is an actual punctured morphism of degree twelve:
its entire origin fibre is precisely \(S\).
Every composite with a deck automorphism of \(X\)
has the same property.

We also record why the target normalizer theorem
holds at every prime in this CM case. The proof of
A.12 applies with \(\mathbf Q(\sqrt{-3})\) in
place of its two listed fields. At a nonsplit prime,
including two and the ramified prime three, an
ordinary good Frobenius and all its positive powers
generate the quadratic CM algebra over
\(\mathbf Q_\ell\). Its eigenvalue quotient has
infinite order, as their ideals have different
support. Thus no rational proper homology line is
stable under an open Galois subgroup, and A.11 applies.
Every split prime is at least seven. The hypotheses
of \cite[Theorem 4.5 and Corollary 4.10]{Ishii2025}
are satisfied: class number one, the full integer
ring, and potentially ordinary reduction at the
split coefficient prime. The index \((2,2)\) obeys
the unit congruence modulo six. Those results make
\(\kappa_{(2,2)}\) nonzero. Nakamura's actual
metabelian formula
\cite[Theorem A and Corollary 4.12]{Nakamura1995}
and the three distinct CM weights isolate the same
middle special derivation \(T\) of A.12, for which
\(T(x)\notin[x,L_4]\) and \(T(y)\notin[y,L_4]\).
The actual lower-inner-term removal, torus projection,
and necessary normalizer equation are exactly those
proved there. They exclude both split eigenlines.
Nonvanishing persists after every finite field
extension. Thus (A.50.9) has the required all-prime
normalizer target, including coefficients two and three.

\paragraph{The rational isotypic factors and the exact span of the maps.}
We give the representation calculation because merely
counting an elliptic quotient would not detect the
correct multiplicity. Let \(F=\mathbf Q(\zeta_3)\).
The nontrivial-central part of \(\mathbf Q[H_3]\)
is \(M_3(F)\). Explicitly its three-dimensional
representation over \(F\) is generated by
\[
 U=\operatorname{diag}(1,\zeta_3,\zeta_3^2),
       \qquad V e_j=e_{j+1}\quad(j\bmod3),
       \qquad UV=\zeta_3VU.                       \tag{A.50.10}
\]
The remaining rational summands are the characters
of \(H_3/\langle c\rangle\). Their representation
on \(H_1(C,\mathbf Q)\) factors through
\(H_1(C/\langle c\rangle,\mathbf Q)=H_1(E,\mathbf Q)\),
on which the quotient group acts by translations and
therefore trivially. Since \(\dim H_1(C)=20\), the
remaining part has rational dimension eighteen,
or three copies of the six-dimensional rational
simple module belonging to \(M_3(F)\).
Rational matrix idempotents therefore give
\(J_C\sim E\times B_3^3\), where
\(\dim B_3=3\) and \(F\) acts on \(B_3\).
These are actual isogeny factors: clear denominators
in the idempotents and take their images in \(J_C\).

The lift (A.50.6) and its \(\tau\)-conjugate can
be chosen so that \(\tau\) sends them to one
another and sends the second to the inverse of the
first. Indeed \(\tau^2\) lies over \([-1]\),
fixes \(x,z\), and conjugates (A.50.6) to its
inverse. In (A.50.10) the induced automorphism is
therefore \(U\mapsto V,V\mapsto U^{-1}\), after
choosing the central generator consistently.
A convenient intertwiner is
\[
 A=\frac1{-\sqrt{-3}}
       (\zeta_3^{-jk})_{0\le j,k\le2}.
                                                               \tag{A.50.11}
\]
Direct multiplication gives
\[
 AUA^{-1}=V,\quad AVA^{-1}=U^{-1},\quad A^4=1,
 \quad\det(t-A)=(t-1)(t^2+1).                   \tag{A.50.12}
\]
On the multiplicity space of dimension three over
\(F\), write \(\tau=A\otimes B_\tau\).
This is possible by the matrix-algebra commutant;
\(B_\tau^4=1\), and its characteristic polynomial
has coefficients in \(F\). Equivalently
\(A^{-1}\tau\) is an actual rational endomorphism
commuting with \(M_3(F)\), hence an endomorphism
of the isogeny factor \(B_3\).

The \(\tau\)-invariant space is precisely
\(q_0^*H_1(E_0,\mathbf Q)\). It has dimension
one over \(F\), because the central automorphism
on this elliptic quotient is multiplication by
\(\zeta_3\). Let \(n_1,n_{-1},k,k\) be the
multiplicities of \(1,-1,i,-i\) in \(B_\tau\).
The last two are equal because \(i\notin F\).
From (A.50.12), the invariant dimension is
\(n_1+2k=1\), while \(n_1+n_{-1}+2k=3\).
Thus
\[
                         n_1=1,\quad n_{-1}=2,\quad k=0.
                                                               \tag{A.50.13}
\]
In particular the invariant vector is a pure tensor:
the one-eigenline of \(A\) tensored with the
one-eigenline of \(B_\tau\). Its \(H_3\)-span
is exactly one copy of the three-dimensional
\(F\)-module, rather than an unspecified larger
isotypic span. The corresponding abelian subvariety
is isogenous to \(E_0^3\), since it is generated
by the images of the actual maps
\(q_0g\), \(g\in H_3\). The complementary
minus-one multiplicity factor has dimension two;
call it \(B\). It inherits the \(F\)-action.
This proves (A.50.2) and the exact dimension of the
homology detected by the punctured maps (A.50.9).
All projectors and maps are defined over \(L\).

\paragraph{A genus-four quotient gives exact Frobenius certificates.}
Fix the lift \(\widetilde T\) in (A.50.6), with
\(x(T)=a\), and set
\(u=z/(x-a)\), \(\eta=Y_a\).
Equations (A.50.4)--(A.50.5) give the actual quotient
\(D=C/\langle\widetilde T\rangle\) as
\[
 \eta^2=\frac{u^9}{27}-a u^6+(3-a^2)u^3+15a-7a^3.
                                                               \tag{A.50.14}
\]
It has genus four. For instance the noncentral
order-three subgroup acts freely on \(C\), since
all inertia in the full cover is central, and
Riemann--Hurwitz gives that genus. Its invariant
space in each three-dimensional representation
(A.50.10) has dimension one. Hence the actual
isogeny decomposition over \(L\) is
\[
                              J_D\sim E\times E_0\times B.
                                                               \tag{A.50.15}
\]
This identification uses the same matrix idempotents
and the same surface \(B\), not an unrelated
surface reconstructed from point counts.

The primes \(61\) and \(313\) split in \(L\):
\(i\) is rational modulo either, and the exact
certificate verifies that all eight nonzero points
of \(E[3]\) are rational over each residue field.
Choose places sending \(a\) to \(4\) and
\(26\), respectively. All curves and maps just
used have good reduction there. Direct exhaustive
counts of (A.50.14) over both the prime field and
its quadratic extension give
\[
\begin{array}{c|r|r|r|r|r}
p&a&\#D(\mathbf F_p)&\#D(\mathbf F_{p^2})
  &\operatorname{tr}(E)&\operatorname{tr}(E_0)\\\hline
61&4&96&3696&-10&-13\\
313&26&312&97872&26&-13
\end{array}                                                     \tag{A.50.16}
\]
The certificate constructs each quadratic field as
\(\mathbf F_p[t]/(t^2-d)\) with a specified
nonsquare \(d\), evaluates the quadratic character
at every argument, and includes the single point
at infinity of the odd-degree hyperelliptic model.
It separately verifies separability of that model
and all curve-map identities.

Subtract the two elliptic traces and their square
traces from (A.50.16). Reciprocity then determines
the two degree-four Frobenius polynomials of the
same surface \(B\):
\[
\begin{aligned}
 P_{61}(t)&=t^4+11t^3+60t^2+671t+3721,\\
 P_{313}(t)&=t^4+11t^3-192t^2+3443t+97969.
\end{aligned}                                                  \tag{A.50.17}
\]
Both have the form
\[
                  t^4+b t^3+(b^2-p)t^2+pb t+p^2,
                         \qquad b=11.              \tag{A.50.18}
\]
If \(\alpha,\bar\alpha\) solve
\(t^2-bt+p=0\), its four roots are
\(\zeta_3\alpha,\zeta_3^2\alpha,
  \zeta_3\bar\alpha,\zeta_3^2\bar\alpha\).
Thus the third-power Frobenius polynomial is
\[
             \bigl(t^2-(b^3-3pb)t+p^3\bigr)^2.
                                                               \tag{A.50.19}
\]
The two ordinary imaginary quadratic fields are
\[
             \mathbf Q(\sqrt{-123}),\qquad
             \mathbf Q(\sqrt{-1131}),               \tag{A.50.20}
\]
which differ. These fields survive every further
positive power: \(p\nmid b\), so the ideals of
\(\alpha\) and \(\bar\alpha\) are supported
at different primes above \(p\). Their ratio is
not a root of unity. Consequently \(\alpha^r\)
is nonreal and generates its same imaginary
quadratic field for every \(r>0\).

\paragraph{Why these counts exclude every CM factor of the surface.}
We spell out the required implication. A surface
with an action by an imaginary quadratic field
\(F\) which has any geometric CM subvariety
must itself be isogenous to the square of a CM
elliptic curve, unless it is geometrically simple
of CM type. In the nonsimple case, a CM elliptic
factor and a non-CM elliptic factor would make the
\(F\)-action project to an impossible unital
embedding \(F\hookrightarrow\mathbf Q\).
Two nonisogenous CM elliptic factors would force
both endomorphism fields to equal \(F\), making
them isogenous after all. The only possibility
is therefore an elliptic square with CM by some
imaginary quadratic field \(F'\).

The geometrically simple CM alternative is impossible
in dimension two with this \(F\)-action. Its
endomorphism field would be a quartic CM field
containing \(F\), hence the compositum of \(F\)
with its real quadratic subfield, a biquadratic
field. Every CM type of a biquadratic CM field is
induced from one of its imaginary quadratic
subfields. The rational Hodge structure is then
two copies of that of a CM elliptic curve, giving
an elliptic-square isogeny after clearing lattice
denominators, as in the Hodge argument of A.3.
This contradicts geometric simplicity.

Suppose, then, that \(B\) had any CM factor.
Over a finite number field it would be isogenous
to \(F_1^2\), with \(F_1\) an elliptic curve
with CM by a fixed field \(F'\). At each of
(61) and (313), pass to a residue extension which
defines this isogeny and all endomorphisms, and
then take a further multiple of three. Formula
(A.50.19) gives ordinary elliptic Frobenius roots
\(\alpha^r,\bar\alpha^r\). The endomorphism
field of an ordinary CM elliptic reduction is
its imaginary quadratic Frobenius field, by
\cite[Theorem V.3.1]{fghSilvermanAEC}; since
it contains \(F'\), it equals \(F'\).
By (A.50.20) and its persistence under powers,
this would force \(F'\) to equal two different
fields. This contradiction proves that \(B\)
has no geometric CM subvariety.

It follows at every coefficient prime that
\(V_\ell(B^3)\) has no one-dimensional
\(\mathbf Q_\ell\)-subrepresentation of any
open absolute Galois subgroup. Indeed such a line
is fixed by that subgroup's closed commutator.
Scale a nonzero vector into its Tate lattice and
reduce modulo \(\ell^n\); this gives torsion
of unbounded order over the maximal abelian
extension of the corresponding finite number
field. \cite[\S1, Corollary 1]{Ruppert1998}
would then give a CM subvariety, which has just
been excluded. No assertion about a conjectural
monodromy group of \(B\) is needed.

\paragraph{All procyclic normalizers.}
Suppose a noncuspidal nontrivial procyclic subgroup
has open arithmetic normalizer. A.11 supplies a
nonzero proper homology vector \(v\) spanning
a rational line stable under an open Galois
subgroup. The actual map \([3]q:X\to U\)
and A.12 make its \(E\)-projection zero.
Every actual map (A.50.9), precomposed with a deck
automorphism, has the CM target recognition proved
above. Each therefore kills the corresponding
proper homology image of \(v\): otherwise the
image procyclic subgroup would be nontrivial,
would retain open arithmetic normalizer, and
would have nonzero proper homology, contrary to
cuspidality on the target.

The exact span calculation (A.50.13) says that these
maps detect all of \(E_0^3\). Thus \(v\) lies
in \(V_\ell(B^3)\), which has no such rational
line. This is a contradiction for every prime,
including three. Notice that A.22 was used only
for the separate full fixed-group assertion at
three, not as a normalizer theorem. Conversely
all nonzero closed cusp inertia subgroups have
open arithmetic normalizer after finite extension
making the cusp rational. This proves the exact
all-prime normalizer classification.

\paragraph{An independent residue-two check.}
After a finite extension the original \(E\) has
good supersingular reduction at two. Its three-torsion
sections remain distinct. Formula (A.50.3), after
removing a constant vertical factor and a finite
extension extracting its cube root, gives a tame
cyclic-three normalization with smooth reduction:
locally at each branch section its order is prime
to three, and elsewhere it is etale. Thus \(C\)
has good reduction. Its order-four automorphism
extends to this model. The special \([-1]\) on
\(E\) has only the origin as a fixed geometric
point, because the elliptic special fibre has
two-rank zero. There is just one point of \(C\)
over that origin. Hence \(\tau^2\), and also
\(\tau\), have exactly this one fixed point;
\(\tau^2\) is nontrivial since it lies over
the nontrivial negation of the elliptic curve.
The special \(\tau\)-action therefore has
order four and just one short orbit, of size one.

Take also a good model of \(E_0\), whose reduction
at two is supersingular. The actual map \(q_0\)
extends to it by the Neron mapping property.
The pullback of a line bundle of degree one has
degree four on every fibre, so the special map is
nonconstant of degree four. It is invariant under
the faithful special \(C_4\)-action. Factoring
through its degree-four function-field quotient
leaves a map of degree one, so that quotient is
exactly the supersingular \(E_0\), without an
inseparable remaining factor. The Deuring--Shafarevich
formula \cite[formula (6), p.~5]{fghGiuliettiKorchmaros}
gives
\[
                  f_C-1=4(0-1)+(4-1)=-1.
                                                               \tag{A.50.21}
\]
Thus \(f_C=0\), and A.46 independently proves
the pro-two normalizer assertion. The nine cusps
are distinct on this good special curve, so this
calculation does not meet the rational joining-path
hypothesis of A.42 and is not promoted to a
fixed-group conclusion at other primes.
\end{proof}
